\documentclass[10pt,a4paper]{amsart}
\usepackage[margin=19mm]{geometry}
\usepackage{amsmath,amssymb,mathtools}
\usepackage[scr=rsfs,cal=euler]{mathalfa}
\usepackage{xfrac}
\usepackage[unicode,hidelinks,bookmarks=false]{hyperref}
\newcommand{\ie}{i.e.\ }
\newcommand{\cf}{cf.\ }
\newcommand{\resp}{respectively\ }
\newcommand{\Subsection}{\subsection}
\providecommand{\nobreakdash}{\nobreak\hbox{-}}
\setlength{\emergencystretch}{2em}
\multlinegap0pt
\usepackage{mathtools}
\usepackage{amssymb}
\usepackage{bm}
\usepackage{algorithm}\usepackage{algpseudocode}
\usepackage{xcolor}
\newtheorem{theorem}{Theorem}
\newtheorem{lemma}{Lemma}
\usepackage{cleveref}
\crefname{theorem}{Theorem}{Theorems}
\crefname{lemma}{Lemma}{Lemmas}
\newcommand{\norm}[1]{\left\| #1 \right\|}
\newcommand{\vect}[1]{\mathbf{#1}}
\title{A Unifying View of Anchoring via Operator-Side Tikhonov Regularization}
\author{Zihao Chen}
\date{Source: arXiv:2605.30905v1 (CC BY 4.0)}
\begin{document}
\maketitle

\noindent\emph{Source and licence.} A Unifying View of Anchoring via Operator-Side Tikhonov Regularization. Version arXiv:2605.30905v1 (CC BY 4.0), distributed by arXiv under the Creative Commons Attribution 4.0 International licence, http://creativecommons.org/licenses/by/4.0/, as recorded in the arXiv licence metadata for this exact version and shown on the abstract page https://arxiv.org/abs/2605.30905v1. Full licence notice: \texttt{licenses/arxiv-2605.30905-CC-BY-4.0.txt}.\par\smallskip

\noindent\textit{Editorial scope.} Counted unit: the article's theorem thm:reg\_forward (source0.tex lines 1133--1147). The article's complete original proof of it is delivered with it (source0.tex lines 1177--1339, one \texttt{proof} environment of 163 lines, which the article places away from the statement). No mathematics is written, repaired or re-derived: the statement and the proof are the article's own lines, in the article's own notation, and no retained line is substituted or re-flowed.  Retained uncounted as the source-local context the proof needs: lines 1105--1111; lines 2129--2149.  Nothing was omitted from any retained span, so the package carries no omission record.  No retained line contains a citation, so the wrapper carries no bibliography and no bibliography entry.  No retained line contains a figure environment, an image-inclusion command or drawing code, so no image is delivered and the package declares no asset dependency.  Editorial formatting only, in addition: an AMS \texttt{amsart} preamble that restates the article's own declarations and macros used by the retained lines, namely the environments \texttt{theorem}, \texttt{lemma} on separate counters, together with the standard packages those lines need.  The retained environments print in this document as Theorem 1, Lemma 1 in order of appearance, whereas the article prints the same items with the numbering of its own full document.  The complete native source of the article as distributed in the arXiv e-print for this exact version is bundled unchanged as \texttt{sources/201700/source0.tex}.
\begin{equation}\label{eq:reg_forward_step}
    \vect{x}_{k+1}
    =
    \vect{x}_k-G_k(\vect{x}_k)
    =
    \vect{x}_k-\eta_kF(\vect{x}_k)-\delta_k(\vect{x}_k-\vect{x}_0).
\end{equation}

\begin{lemma}\label{lem:recurrence}
Let \(c,c_1,\alpha,\beta,\gamma>0\) and \(R\ge0\) with $1 < \gamma \alpha \le \beta$ and define
\[
   \delta_n \;=\; \frac{\alpha}{\,n+\beta\,}, 
   \qquad n\ge0.
\]
Assume the sequence \((e_n)_{n\ge 1}\) satisfies $e_1 \le c_1\,R$ and
\begin{equation} \label{eq:rec}
   e_{n+1}  \le \bigl(1-\gamma\delta_n\bigr)
         \Bigl(e_{n}+c\,(\delta_{n-1}-\delta_n)\,R\Bigr),
         \qquad n\ge 1. 
\end{equation}
Then for every \(n\ge 1\),
\[
   e_n \le \frac{M}{n + \beta - 1}, \quad \text{where }
        M := R\max\!\bigl\{
            \tfrac{c\,\alpha}{\gamma\alpha-1},
            \;c_1\beta
         \bigr\}.
\]
\end{lemma}

\begin{theorem}[Tikhonov-Regularized Forward Residual Bound]\label{thm:reg_forward}
Let \(F:\mathcal H\to\mathcal H\) be a monotone, \(L\)-Lipschitz operator with
\(L>0\) and with a zero \(\vect{x}^*\in\mathcal H\). Let
\(R\coloneqq\norm{\vect{x}_0-\vect{x}^*}\). Consider the iteration
\eqref{eq:reg_forward_step} with parameters satisfying
\(2\eta_kL=\sqrt{\delta_k}\) and \(\delta_k=4/(k+4)\) for \(k\ge0\). Then, for
every \(k\ge1\),
\[
    r_F(\vect{x}_k)
    =
    \norm{F(\vect{x}_k)}
    \le
    \frac{24LR}{\sqrt{k+3}}.
\]
\end{theorem}

\begin{proof}[Proof of \cref{thm:reg_forward}]
We carry out the three steps in turn.

\paragraph{Step 1: The regularized update is a contraction.}
Consider the map \(S_k(\vect{x})\coloneqq\vect{x}-G_k(\vect{x})\). For any
\(\vect{x},\vect{y}\in\mathcal H\),
\[
    S_k(\vect{x})-S_k(\vect{y})
    =
    (1-\delta_k)(\vect{x}-\vect{y})
    -\eta_k\bigl(F(\vect{x})-F(\vect{y})\bigr).
\]
Squaring and using monotonicity of \(F\) to drop the cross term gives
\begin{align*}
    \norm{S_k(\vect{x})-S_k(\vect{y})}^2
    &\le
    (1-\delta_k)^2\norm{\vect{x}-\vect{y}}^2
    +\eta_k^2\norm{F(\vect{x})-F(\vect{y})}^2 \\
    &\le
    \bigl((1-\delta_k)^2+\eta_k^2L^2\bigr)
    \norm{\vect{x}-\vect{y}}^2.
\end{align*}
Substituting \(\eta_k^2L^2=\delta_k/4\) and using
\(\delta_k\in(0,1]\),
\[
    (1-\delta_k)^2+\frac{\delta_k}{4}
    =
    1-2\delta_k+\delta_k^2+\frac{\delta_k}{4}
    \le
    \left(1-\frac{\delta_k}{2}\right)^2.
\]
Thus \(S_k\) is a \((1-\delta_k/2)\)-contraction.

\paragraph{Step 2: The trajectory is bounded.}
Since \(F(\vect{x}^*)=\vect{0}\),
\[
    G_k(\vect{x}^*)=\delta_k(\vect{x}^*-\vect{x}_0).
\]
Therefore,
\[
    \norm{\vect{x}_{k+1}-\bigl(\vect{x}^*-G_k(\vect{x}^*)\bigr)}
    =
    \norm{S_k(\vect{x}_k)-S_k(\vect{x}^*)}
    \le
    \left(1-\frac{\delta_k}{2}\right)\norm{\vect{x}_k-\vect{x}^*}.
\]
The triangle inequality gives
\[
    \norm{\vect{x}_{k+1}-\vect{x}^*}
    \le
    \left(1-\frac{\delta_k}{2}\right)\norm{\vect{x}_k-\vect{x}^*}
    +\delta_k R.
\]
Induction yields \(\norm{\vect{x}_k-\vect{x}^*}\le2R\) for all \(k\ge0\): the
base case is immediate, and if \(\norm{\vect{x}_k-\vect{x}^*}\le2R\), then
\[
    \norm{\vect{x}_{k+1}-\vect{x}^*}
    \le
    \left(1-\frac{\delta_k}{2}\right)2R+\delta_kR
    =
    2R.
\]
Consequently,
\[
    \norm{\vect{x}_k-\vect{x}_0}
    \le
    \norm{\vect{x}_k-\vect{x}^*}+\norm{\vect{x}^*-\vect{x}_0}
    \le
    3R.
\]

\paragraph{Step 3: Recurrence and final bound.}
Using \(\vect{x}_{k+1}=\vect{x}_k-G_k(\vect{x}_k)\) and the contraction of
\(S_k\),
\begin{align*}
    \norm{G_k(\vect{x}_{k+1})}
    &=
    \norm{S_k(\vect{x}_{k+1})-S_k(\vect{x}_k)}
    \qquad(\text{since }\vect{x}_{k+1}=S_k(\vect{x}_k)) \\
    &\le
    \left(1-\frac{\delta_k}{2}\right)\norm{\vect{x}_{k+1}-\vect{x}_k}
    =
    \left(1-\frac{\delta_k}{2}\right)\norm{G_k(\vect{x}_k)}.
\end{align*}
We next compare \(G_k(\vect{x}_k)\) to
\(E_k=\norm{G_{k-1}(\vect{x}_k)}\). From the definition of \(G_k\),
\[
    G_k(\vect{x}_k)
    =
    \frac{\eta_k}{\eta_{k-1}}G_{k-1}(\vect{x}_k)
    +
    \left(\delta_k-\frac{\eta_k}{\eta_{k-1}}\delta_{k-1}\right)
    (\vect{x}_k-\vect{x}_0).
\]
Since \(\delta_k\le\delta_{k-1}\) and
\(\eta_k/\eta_{k-1}=\sqrt{\delta_k/\delta_{k-1}}\le1\),
\[
    \left|\delta_k-\frac{\eta_k}{\eta_{k-1}}\delta_{k-1}\right|
    =
    \sqrt{\delta_k\delta_{k-1}}-\delta_k
    \le
    \delta_{k-1}-\delta_k.
\]
Together with \(\norm{\vect{x}_k-\vect{x}_0}\le3R\), this gives
\[
    \norm{G_k(\vect{x}_k)}
    \le
    \frac{\eta_k}{\eta_{k-1}}E_k
    +(\delta_{k-1}-\delta_k)\norm{\vect{x}_k-\vect{x}_0}
    \le
    E_k+3(\delta_{k-1}-\delta_k)R.
\]
Setting \(E_{k+1}=\norm{G_k(\vect{x}_{k+1})}\), we obtain, for \(k\ge1\),
\[
    E_{k+1}
    \le
    \left(1-\frac{\delta_k}{2}\right)
    \bigl(E_k+3(\delta_{k-1}-\delta_k)R\bigr).
\]
The base case follows from the contraction of \(S_0\), the identity
\(G_0(\vect{x}_0)=\eta_0F(\vect{x}_0)\), and Lipschitzness around
\(\vect{x}^*\):
\[
    E_1
    =
    \norm{G_0(\vect{x}_1)}
    \le
    \left(1-\frac{\delta_0}{2}\right)\norm{\eta_0F(\vect{x}_0)}
    \le
    \frac{R}{4}.
\]
Here \(\delta_0=1\), \(\eta_0=1/(2L)\), and
\(\norm{F(\vect{x}_0)}=\norm{F(\vect{x}_0)-F(\vect{x}^*)}\le LR\).
Applying \cref{lem:recurrence} with
\[
    \alpha=\beta=4,\qquad
    \gamma=\frac12,\qquad
    c=3,\qquad
    c_1=\frac14
\]
yields
\[
    E_k\le\frac{12R}{k+3}.
\]
Finally, since
\[
    G_{k-1}(\vect{x}_k)
    =
    \eta_{k-1}F(\vect{x}_k)
    +\delta_{k-1}(\vect{x}_k-\vect{x}_0),
\]
we have
\[
    \norm{F(\vect{x}_k)}
    \le
    \frac{E_k+3\delta_{k-1}R}{\eta_{k-1}}
    \le
    \frac{24LR}{\sqrt{k+3}},
\]
where the last inequality uses
\(\delta_{k-1}=4/(k+3)\), \(\eta_{k-1}=1/(L\sqrt{k+3})\), and
\(E_k\le12R/(k+3)\).
\end{proof}

\end{document}
